% TEMPLATE-MILC-v3.tex - plantilla maestra UNICA de las guias MILC v3.
%
% La usan las cuatro familias de guias, cada una con su driver:
%   · Grados 6-11          -> scripts/build-guias-g11.py
%   · Bebras               -> scripts/build-guias-bebras.py
%   · Semillero            -> scripts/build-guias-semillero.py
%   · Territorio interior  -> scripts/build-guias-territorio-interior.py
%
% Antes habia tres copias casi identicas (~400 lineas iguales, 6 distintas):
% cada mejora habia que aplicarla tres veces, y en la practica no se hacia
% (el motor de recursos vivia solo en dos de ellas). Lo que varia entre
% familias esta parametrizado en 6 placeholders que cada driver llena
% (se nombran aqui SIN su sintaxis real: el builder reemplaza por texto plano,
% tambien dentro de los comentarios, y un valor multilinea rompe el preambulo):
%
%   HEADER_IZQ        encabezado izquierdo de pagina
%   HEADER_DER        encabezado derecho de pagina
%   PORTADA_ETIQUETA  rotulo sobre el numero grande (GRADO/MOMENTO/MODULO)
%   PORTADA_SERIE     linea de serie de la portada
%   PORTADA_ID        identificador de la guia en la portada
%   PORTADA_CIRCULO   el circulito de grado (°), o vacio si no aplica
%   PDF_TITULO        titulo en texto plano (sin \textbf ni comillas TeX) para
%                     los metadatos del PDF (/Title); lo produce lib_xelatex.texto_pdf
%
% Composicion (auditoria 2026-09-03): las bandas de estacion piden espacio
% (\Needspace*) y prohiben el salto justo despues (\nopagebreak), asi no quedan
% huerfanas al pie; las cajas largas (softbox, drawbox, infoband, puentebox) son
% `breakable`, asi una caja de media pagina ya no arrastra todo a la siguiente
% dejando la anterior al 40 %. Todo texto sobre mostaza o verde va en milcNegro
% (blanco sobre #E5B400 da 1,93:1 y sobre #58A923 2,95:1; ninguno pasa WCAG AA).
% El resto de claves las documenta content/guias/_SCHEMA.md. El script valida
% que no queden placeholders sin reemplazar antes de escribir el archivo final.

% Etiquetado PDF (latex-lab): /Lang, estructura y texto alternativo de las
% figuras, para lectores de pantalla. Debe ir ANTES de \documentclass.
\DocumentMetadata{lang=es-CO,tagging=on}
\documentclass[11pt,letterpaper]{article}
% headheight=24pt: la cabecera derecha lleva dos lineas (identidad de la guia
% y estacion actual). headsep se acorta en la misma medida para que el bloque
% de texto quede exactamente donde estaba con headheight=12pt.
\usepackage[margin=2.0cm,headheight=24pt,headsep=13pt]{geometry}
\usepackage[table]{xcolor}
\usepackage{fontspec}
\usepackage[spanish]{babel}
\usepackage{tikz}
% Librerías TikZ para los diagramas de las guías (bloque `recursos:`):
%  · babel  → babel en español vuelve activos caracteres como «>», y TikZ los usa
%             en sus opciones (>=stealth, ->). Sin esta librería, cualquier
%             diagrama con flechas revienta con «\language@active@arg> has an
%             extra }». Debe ir ANTES de que se dibuje nada.
%  · positioning        → sintaxis «below left=2pt and 3pt».
%  · decorations.pathreplacing → llaves ({brace}) para señalar rangos.
\usetikzlibrary{babel,positioning,decorations.pathreplacing}
\usepackage{graphicx}
% Circuitos: el trabajo de electrónica se hace hoy con micro:bit y sus sensores
% integrados (MakeCode, sin cableado), así que no se necesitan esquemáticos.
% Si algún día entran montajes con protoboard, basta descomentar esta línea:
% los diagramas .tex/.tikz declarados en `recursos:` ya se incrustan solos.
% \usepackage{circuitikz}
\usepackage[most]{tcolorbox}
\usepackage{tabularx}
\usepackage{array}
\usepackage{fancyhdr}
\usepackage{titlesec}
\usepackage[document]{ragged2e}  % bandera a la derecha en todo el cuerpo
\usepackage{enumitem}
\usepackage{microtype}
\usepackage{etoolbox}
\usepackage{needspace}
% hyperref al final del preambulo: metadatos del PDF (titulo, autor, idioma,
% familia), marcadores por estacion (\pdfbookmark) y enlaces sin recuadro.
\usepackage[hidelinks,
  pdftitle={Empatía no es lástima: acompañar es un oficio, no un sentimiento},
  pdfauthor={PhD. Álvaro Cárdenas Orozco},
  pdfsubject={Territorio interior · Educación socioemocional · Grado 7° · Momento 3 · MILC v3},
  pdflang={es-CO},
  bookmarksopen=true,bookmarksnumbered=false]{hyperref}

% Helvetica Neue no trae la flecha U+2192, asi que cada «→» escrito literalmente
% en el YAML se compilaba a NADA, en silencio (462 flechas en 61 guias: rutas del
% tipo «paleta → tipografia → lockup» salian rotas). Mapeamos el caracter a su
% version matematica, que si tiene glifo. Asi el autor puede seguir escribiendo
% «→» comodamente en el YAML.
\usepackage{newunicodechar}
\newunicodechar{→}{\ensuremath{\rightarrow}}
% Mismo problema con los simbolos de marca y vinetas: el «✓» de «una palomita ✓»
% salia como cuadrito vacio en el PDF de todas las guias que lo usaban.
\usepackage{pifont}
\newunicodechar{✓}{\ding{51}}
\newunicodechar{✗}{\ding{55}}
\newunicodechar{✔}{\ding{52}}
\newunicodechar{•}{\textbullet}
\newunicodechar{≥}{\ensuremath{\geq}}
\newunicodechar{≤}{\ensuremath{\leq}}

\setmainfont{Helvetica Neue}
\setsansfont{Helvetica Neue}
\setlength{\parindent}{0pt}
\setlength{\parskip}{6pt}
% Sin particion de palabras: en 6.º-7.º una palabra cortada por guion al final
% de linea («Comuni-cacion») frena la lectura mucho mas que una linea corta.
\hyphenpenalty=10000
\exhyphenpenalty=10000
\pretolerance=2000
\tolerance=3000
\setlength{\emergencystretch}{3em}
\linespread{1.10}
\renewcommand{\arraystretch}{1.25}
\setlist{leftmargin=5mm,itemsep=1mm,topsep=1mm}
\newcolumntype{Y}{>{\RaggedRight\arraybackslash}X}

\definecolor{milcMagenta}{HTML}{E6007E}
\definecolor{milcVino}{HTML}{5A0038}
\definecolor{milcTurquesa}{HTML}{0093A5}
\definecolor{milcVerde}{HTML}{58A923}
\definecolor{milcMostaza}{HTML}{E5B400}
\definecolor{milcOcre}{HTML}{B86600}
\definecolor{milcNegro}{HTML}{241816}
\definecolor{milcGris}{HTML}{F7F7F4}
\definecolor{milcLinea}{HTML}{E8E2DD}
\definecolor{milcGradePrimary}{HTML}{0D9488}
\definecolor{milcGradeDark}{HTML}{115E59}
\definecolor{milcGradeSoft}{HTML}{CCFBF1}
\definecolor{milcGradeAccent}{HTML}{CFE7FF}
\definecolor{milcGradeText}{HTML}{FFFFFF}
\color{milcNegro}

\pagestyle{fancy}
\fancyhf{}
% Cabecera de dos lineas: identidad de la guia arriba y, a la derecha y en
% gris, la estacion en curso (\markright desde \milcestacion). Antes de la
% primera estacion la segunda linea va vacia; el \strut mantiene la altura
% para que la linea de arriba no baile entre paginas.
\lhead{\small Territorio interior · Educación socioemocional\\[-1pt]{\footnotesize\strut}}
\rhead{\small Grado 7° · Momento 3 · MILC v3\\[-1pt]{\footnotesize\strut\color{milcNegro!65}\rightmark}}
\cfoot{\small\thepage}
\renewcommand{\headrulewidth}{0pt}

% Sobre mostaza (#E5B400) y verde (#58A923) el texto va en milcNegro; sobre el
% resto de la paleta, en blanco. Se usa dentro de fonttitle/fontupper, donde un
% \color posterior manda sobre coltitle.
\newcommand{\milctintasobre}[1]{%
  \ifboolexpr{test{\ifstrequal{#1}{milcMostaza}} or test{\ifstrequal{#1}{milcVerde}}}%
    {\color{milcNegro}}{\color{white}}}

\newtcolorbox{titlebox}[1]{
  enhanced,colback=#1,colframe=#1,arc=7mm,boxrule=0pt,
  left=7mm,right=7mm,top=3mm,bottom=3mm,
  before skip=8pt,after skip=8pt,
  fontupper=\sffamily\bfseries\Large\color{white}
}

\newtcolorbox{softbox}[3]{
  enhanced,breakable,pad at break=3mm,
  colback=#2,colframe=#1,borderline west={4pt}{0pt}{#1},
  arc=2mm,boxrule=.4pt,left=4mm,right=4mm,top=3mm,bottom=3mm,
  before skip=6pt,after skip=8pt,title={#3},coltitle=white,
  colbacktitle=#1,fonttitle=\sffamily\bfseries\milctintasobre{#1},fontupper=\RaggedRight,
  attach boxed title to top left={xshift=3mm,yshift=-2mm},
  boxed title style={arc=2mm, boxrule=0pt}
}

\newtcolorbox{drawbox}[2]{
  enhanced,breakable,pad at break=4mm,
  colback=white,colframe=#1,arc=4mm,boxrule=1pt,
  left=5mm,right=5mm,top=4mm,bottom=4mm,before skip=8pt,after skip=8pt,
  title={#2},coltitle=white,colbacktitle=#1,fonttitle=\sffamily\bfseries\milctintasobre{#1},
  fontupper=\RaggedRight,
  attach boxed title to top left={xshift=4mm,yshift=-2mm},
  boxed title style={arc=3mm, boxrule=0pt}
}

\newtcolorbox{infoband}[1]{
  enhanced,breakable,pad at break=2mm,colback=#1!8,colframe=#1!50,arc=2mm,boxrule=.5pt,
  left=4mm,right=4mm,top=2mm,bottom=2mm,before skip=4pt,after skip=4pt,
  fontupper=\small\RaggedRight
}

% Puente narrativo entre secciones (contrato editorial Conéctate).
% Da continuidad: "Acabas de X. Ahora vas a Y porque Z."
% Visualmente: banda izquierda en color del grado, texto en itálica pequeña.
\newtcolorbox{puentebox}{
  enhanced,breakable,pad at break=2mm,colback=milcGradeSoft,colframe=milcGradeDark,
  borderline west={3pt}{0pt}{milcGradeDark},
  arc=0mm,boxrule=0pt,left=5mm,right=4mm,top=2.5mm,bottom=2.5mm,
  before skip=4pt,after skip=8pt,
  fontupper=\itshape\small\color{milcNegro}\RaggedRight
}
% La flecha va en modo matematico a proposito: Helvetica Neue NO tiene el glifo
% U+2192, asi que \textrightarrow se compilaba a NADA («Missing character: There
% is no → in font Helvetica Neue») y el puente salia sin flecha. En modo
% matematico la flecha viene de la fuente matematica y si se dibuja.
\newcommand{\puente}[1]{%
  \begin{puentebox}%
    \textcolor{milcGradeDark}{$\rightarrow$}\hspace{2mm}#1%
  \end{puentebox}%
}

\newcommand{\linead}[1]{\noindent\color{milcLinea}\rule{#1}{0.5pt}\color{milcNegro}}
\newcommand{\respuesta}[1]{\par\vspace{2mm}\foreach \n in {1,...,#1}{\noindent\color{milcLinea}\rule{\linewidth}{0.5pt}\par\vspace{5mm}}\color{milcNegro}}
\newcommand{\checkbox}{\textcolor{milcGradeDark}{\fbox{\phantom{x}}}\hspace{5pt}}

% ─── Listas aireadas para las guías socioemocionales ────────────────────
% Componer las verificaciones como «\checkbox texto\\» deja el texto que salta
% de línea debajo del cuadrito y apelmaza el bloque. Estos entornos alinean la
% continuación y airean los ítems. Los usa Territorio Interior; cualquier
% familia puede hacerlo.
\newlist{checklista}{itemize}{1}
\setlist[checklista]{
  label=\textcolor{milcGradeDark}{\fbox{\phantom{x}}},
  leftmargin=9mm, labelsep=3.2mm, itemsep=2.4mm,
  topsep=2.5mm, parsep=0pt, partopsep=0pt, before=\RaggedRight
}
\newlist{pasoslista}{enumerate}{1}
\setlist[pasoslista]{
  label=\textcolor{milcGradeDark}{\sffamily\bfseries\arabic*.},
  leftmargin=8mm, labelsep=3mm, itemsep=2.8mm,
  topsep=1mm, parsep=0pt, partopsep=0pt, before=\RaggedRight
}

\newcommand{\phasecard}[4]{
\begin{tcolorbox}[enhanced,colback=#3,colframe=#2,arc=3mm,boxrule=.5pt,height=3.05cm,valign=center,halign=center,left=1.5mm,right=1.5mm,top=2mm,bottom=2mm]
{\sffamily\bfseries\fontsize{22}{24}\selectfont\textcolor{#2}{#1}}\par
{\sffamily\bfseries\small #4}
\end{tcolorbox}
}

% ── Piezas del rediseno 2026 (legibilidad 6.º-11.º) ───────────────────
% Banda de estacion: reemplaza el titlebox macizo. Titulo a la izquierda,
% dato practico (tiempo/modalidad) a la derecha, en una sola linea.
% Nunca huerfana: pide 9 lineas libres (o salta de pagina rellenando con
% \vfill) y prohibe el salto justo despues de la banda. Nueve y no seis:
% la caja partible que sigue necesita ~80pt (titulo adosado, relleno y dos
% lineas) para arrancar en la misma pagina; con menos, tcolorbox la manda
% entera a la siguiente y la banda se queda sola. El penalty 10000 va
% en `after`, ANTES del espacio posterior: un `after skip` (aunque sea 0pt)
% es un \vskip tras la caja, y ese glue es un punto de corte legal que un
% \nopagebreak escrito despues de \end{tcolorbox} ya no cierra. Cada banda
% deja marcador PDF y actualiza la cabecera derecha.
\newcounter{milcestacion}
\newcommand{\milcestacion}[2]{%
\Needspace*{9\baselineskip}%
\stepcounter{milcestacion}%
\pdfbookmark[1]{#2}{milcestacion\themilcestacion}%
\markright{#2}%
\begin{tcolorbox}[enhanced,colback=#1,colframe=#1,arc=2mm,boxrule=0pt,
  left=5mm,right=5mm,top=2.6mm,bottom=2.6mm,before skip=11pt,
  after={\par\nopagebreak\vskip7pt}]
{\sffamily\bfseries\fontsize{15}{18}\selectfont\milctintasobre{#1}#2}
\end{tcolorbox}}

% Tarjeta de paso numerada: sustituye las tablas «Pilar / Aplicacion», que
% eran formato de consulta docente, no de estudiante. El numero va en un
% circulo sobre el borde para que la secuencia se lea de un vistazo.
\newcommand{\milcpaso}[3]{%
\Needspace{4\baselineskip}%
\begin{tcolorbox}[enhanced,colback=white,colframe=milcLinea,arc=1.5mm,boxrule=.9pt,
  left=14mm,right=4mm,top=3mm,bottom=3mm,before skip=4pt,after skip=4pt,
  fontupper=\RaggedRight,
  overlay={\node[anchor=north west,circle,fill=#2,text=white,inner sep=0pt,
    minimum size=7.4mm,font=\sffamily\bfseries\small]
    at ([xshift=3.6mm,yshift=-3.2mm]frame.north west) {#1};}]
#3
\end{tcolorbox}}

% Casilla marcable de tamano util con lapiz (la anterior era \fbox{\phantom{x}},
% de alto variable segun el texto que la seguia).
\newcommand{\milcmarca}{\framebox[3.8mm]{\rule{0pt}{3.4mm}}}

% Fila marcable de lista suelta (checks de la estacion 1).
\newcommand{\milccheckrow}[1]{%
\par\vspace{1.8mm}\noindent
\makebox[6.4mm][l]{\raisebox{-.55mm}{\textcolor{milcNegro}{\milcmarca}}}%
\parbox[t]{\dimexpr\linewidth-6.4mm\relax}{\RaggedRight #1}\par}

% Celda marcable dentro de la rubrica (tabla de autoevaluacion). Misma
% sangria francesa, pero pensada para vivir dentro de una columna X.
\newcommand{\milcopcion}[1]{%
\makebox[5.2mm][l]{\raisebox{-.55mm}{\textcolor{milcNegro}{\milcmarca}}}%
\parbox[t]{\dimexpr\linewidth-5.2mm\relax}{\RaggedRight #1}}

% ── Recursos visuales de la guía ──────────────────────────────────────
% Declarados en el bloque `recursos:` del YAML y emitidos por el builder.
% \guiaFigura   → imagen rasterizada o PDF (foto, carta celeste, montaje).
% \guiaDiagrama → fuente TikZ/circuitikz (.tex) incrustada como vector:
%                 es la vía para circuitos y esquemáticos editables.
% [#1] = texto alternativo (alt) · #2 = ruta relativa al .tex · #3 = pie
% El alt llega al PDF etiquetado (/Alt) y lo lee el lector de pantalla.
\newcommand{\guiaFigura}[3][]{%
\begin{center}
\includegraphics[width=0.88\linewidth,height=0.38\textheight,keepaspectratio,alt={#1}]{#2}\\[1.5mm]
{\footnotesize\itshape\textcolor{milcNegro!75}{#3}}
\end{center}
}

\newcommand{\guiaDiagrama}[2]{%
\begin{center}
\input{#1}\\[1.5mm]
{\footnotesize\itshape\textcolor{milcNegro!75}{#2}}
\end{center}
}

\begin{document}
\thispagestyle{empty}

% ── Portada ───────────────────────────────────────────────────────────
% Antes: un rectangulo de color a pagina completa. Se veia bien en pantalla,
% pero en la fotocopiadora del colegio se comia el toner y salia gris sucio.
% Ahora la banda ocupa el tercio superior y el resto va en blanco.
\begin{tikzpicture}[remember picture,overlay]
  \path[fill=milcGradePrimary]
    (current page.north west) rectangle ([yshift=-10.6cm]current page.north east);

  \node[anchor=north west,text=milcGradeText] at ([xshift=2.0cm,yshift=-1.55cm]current page.north west)
    {\sffamily\bfseries\fontsize{12}{14}\selectfont MOMENTO};
  \node[anchor=north west,text=milcGradeText,opacity=.85] at ([xshift=2.0cm,yshift=-2.35cm]current page.north west)
    {\sffamily\bfseries\fontsize{8.5}{10}\selectfont TERRITORIO INTERIOR · SOCIOEMOCIONAL};
  \node[anchor=north east,text=milcGradeText,opacity=.85,align=right] at ([xshift=-2.0cm,yshift=-1.55cm]current page.north east)
    {\sffamily\bfseries\fontsize{9}{12}\selectfont I.E. Sor María Juliana\\Cartago, Valle del Cauca};

  \node[anchor=west,text=milcGradeText] (gradonum) at ([xshift=1.85cm,yshift=-7.1cm]current page.north west)
    {\sffamily\bfseries\fontsize{112}{112}\selectfont 3};
  % El circulito de grado (°) solo aplica a las guias de grado («11°»); en
  % Bebras y Semillero el numero grande es un momento/modulo, no un grado.
  % Se posiciona relativo al nodo `gradonum`, no en coordenadas absolutas.
  

  \node[anchor=west,text=milcGradeText,text width=11.4cm,align=left] at ([xshift=1.5cm,yshift=0.5cm]gradonum.east)
    {
     {\sffamily\bfseries\fontsize{9}{11}\selectfont\MakeUppercase{Territorio interior · Socioemocional}}\\[3mm]
     {\sffamily\bfseries\fontsize{21}{25}\selectfont\setlength{\baselineskip}{27pt}Empatía no es lástima\\[4pt]acompañar\\[4pt]es un oficio}\\[3.5mm]
     {\sffamily\bfseries\fontsize{15}{18}\selectfont Grado 7° · Momento 3}
    };
\end{tikzpicture}

\vspace*{9.4cm}
\vfill

{\sffamily\bfseries\fontsize{8}{10}\selectfont\textcolor{milcGradeDark}{LO QUE TE LLEVAS HOY}}

\begin{tcolorbox}[enhanced,colback=white,colframe=milcNegro,arc=2mm,boxrule=1.6pt,
  left=5mm,right=5mm,top=4mm,bottom=4mm,before skip=3pt,after skip=12pt,
  fontupper=\RaggedRight\fontsize{11}{15.5}\selectfont]
Tu guía de acompañamiento: tres situaciones traducidas a la perspectiva del otro, la lista de lo que sí y lo que no se dice, y un acompañamiento real hecho esta semana —con lo que hiciste, no con lo que sentiste.
\end{tcolorbox}

\begin{tcolorbox}[enhanced,colback=milcGris,colframe=milcGris,arc=2mm,boxrule=0pt,
  left=5mm,right=5mm,top=3.5mm,bottom=3.5mm,after skip=12pt]
{\sffamily\bfseries\fontsize{8}{10}\selectfont\textcolor{milcNegro!55}{ESTA GUÍA ES DE}}\\[3mm]
\begin{tabularx}{\linewidth}{X p{3.2cm} p{3.2cm}}
\linead{\linewidth} & \linead{3.0cm} & \linead{3.0cm}\\[-1mm]
{\fontsize{8.5}{10}\selectfont\textcolor{milcNegro!55}{Nombre completo}} &
{\fontsize{8.5}{10}\selectfont\textcolor{milcNegro!55}{Grupo}} &
{\fontsize{8.5}{10}\selectfont\textcolor{milcNegro!55}{Fecha}}\\
\end{tabularx}
\end{tcolorbox}

\vfill

\noindent\textcolor{milcLinea}{\rule{\linewidth}{1pt}}\\[2mm]
{\sffamily\bfseries\fontsize{9.5}{12}\selectfont Autor: Dr. Álvaro Cárdenas Orozco}\\[1mm]
{\sffamily\fontsize{8.5}{11}\selectfont\textcolor{milcNegro!55}{Metodología MILC · Apertura ancestral · Empatía y acompañamiento · Triángulo Dussel-Estoicismo-Floridi}}

\newpage

\section*{Apertura: Empatía no es lástima: acompañar es un oficio, no un sentimiento}

\begin{softbox}{milcGradeDark}{milcGradeSoft}{Saber ancestral}
Cuando alguien muere en el Chocó y en buena parte del Pacífico, \textbf{no se deja solo a nadie: se canta}. Son los \textbf{alabaos}, cantos fúnebres a capela ---sin instrumentos--- que se cantan durante el velorio y el novenario, con una \textbf{voz líder} que empieza y \textbf{un coro de mujeres que responde}. Las \textbf{cantadoras} son especialistas: se sabe a quién llamar, se sabe qué se canta según quién murió ---si es un niño, el ritual es otro, el gualí o chigualo--- y se sabe cuándo termina, porque el último día, mientras ellas cantan, se desmonta el altar y se apagan las velas. Los alabaos son \textbf{Patrimonio Inmaterial de la Nación desde 2014}, y traen adentro dos historias mezcladas: la del catolicismo popular que llegó con la colonización y la de los repertorios espirituales africanos que llegaron con las personas esclavizadas. \textbf{Fíjate en tres cosas que hacen las cantadoras:} llegan y \emph{se quedan} ---noches enteras---; \textbf{no consuelan con frases}, cantan; y \textbf{nadie canta solo}: una empieza y las otras responden. \emph{Eso es acompañar: un oficio con horas, con técnica y con gente. No un sentimiento que se tiene desde lejos.}
\end{softbox}

\begin{softbox}{milcTurquesa}{milcGris}{Saber contemporáneo}
En español usamos una sola palabra ---«empatía»--- para tres cosas que la psicología separa hace cuarenta años. Mark Davis construyó en 1983 el instrumento más usado para medirla y encontró que no es una sino \textbf{varias capacidades distintas} (Davis, 1983). Tres importan aquí. La \textbf{toma de perspectiva}: la tendencia a adoptar el punto de vista del otro, entender \emph{desde dónde} está mirando ---y es más cabeza que corazón---. La \textbf{preocupación empática}: sentir algo por lo que le pasa al otro, que es lo que empuja a hacer algo. Y el \textbf{malestar personal}: la angustia \emph{propia} ante el dolor ajeno. \textbf{Esa tercera es la trampa.} Cuando alguien está mal y tú te angustias tanto que terminas necesitando que te consuelen a ti, o que huyes porque no lo soportas, no estás acompañando: \textbf{estás gestionando tu propio malestar}. \emph{Y la lástima es todavía otra cosa:} sentir pena por alguien mirándolo desde arriba ---«pobrecito»---, que consuela al que la siente y humilla al que la recibe.
\end{softbox}

\begin{softbox}{milcMostaza}{milcGris}{Pregunta-puente}
\textit{Las cantadoras no llegan a decirle a la familia «qué pena, ánimo»: llegan, se sientan y cantan hasta el amanecer. \textbf{¿Cuándo fue la última vez que acompañaste a alguien con algo que hiciste}\ldots y no solo con algo que dijiste?}
\end{softbox}

\begin{softbox}{milcVerde}{milcGris}{Saber y saber hacer}
Al terminar podrás: (1) \textbf{identificar} la diferencia entre toma de perspectiva, preocupación empática, malestar personal y lástima, en casos tuyos; (2) \textbf{explicar} por qué la lástima humilla y qué la distingue del acompañamiento; (3) \textbf{aplicar} la toma de perspectiva contando una situación desde el punto de vista del otro sin descalificarlo; (4) \textbf{crear} tu guía de acompañamiento y usarla con una persona real esta semana.
\end{softbox}



\small
\begin{tabularx}{\textwidth}{>{\bfseries}p{3.0cm}Y}
\rowcolor{milcGradePrimary!12}Autor & Dr. Álvaro Cárdenas Orozco\\
DBA & Comprende los fenómenos que afectan a su territorio, regula sus emociones ante ellos y acompaña a otros con empatía (transversal 6°--11°, articulado con la política «Escuela, territorio de vida» del MEN, Resolución 006519 de 2025, y la Ley 1620 de 2013)\\
Referentes & Primera ayuda psicológica (OMS, 2011/2012) · Directrices IASC (2007) · USGS y Servicio Geológico Colombiano · Pueblo nasa y la minga (CRIC) · Dussel · Estoicismo · Floridi\\
Producto final & Tu guía de acompañamiento: tres situaciones traducidas a la perspectiva del otro, la lista de lo que sí y lo que no se dice, y un acompañamiento real hecho esta semana —con lo que hiciste, no con lo que sentiste.\\
\end{tabularx}
\normalsize

\puente{Ya trabajaste escuchar y los celos. Hoy vas un paso más difícil: \textbf{estar al lado de alguien que la está pasando mal} sin arreglarlo, sin huir y sin ponerte por encima.}

\milcestacion{milcGradeDark}{Ruta MILC: así vamos a aprender}
\begin{tabularx}{\textwidth}{>{\centering\arraybackslash}X>{\centering\arraybackslash}X>{\centering\arraybackslash}X>{\centering\arraybackslash}X}
\phasecard{1}{milcVerde}{milcGris}{Escucha\\[-1mm]\small Observo, escucho y pregunto.} &
\phasecard{2}{milcTurquesa}{milcGris}{Sistematización\\[-1mm]\small Distingo y acompaño.} &
\phasecard{3}{milcMagenta}{milcGris}{Praxis\\[-1mm]\small Construyo y aplico.} &
\phasecard{4}{milcVino}{milcGris}{Evaluación\\[-1mm]\small Reflexión liberadora.}
\end{tabularx}


\puente{Primero, tus propios casos: las veces que te acompañaron bien y las veces que te acompañaron mal. Ahí está todo lo que hay que aprender.}

\milcestacion{milcVerde}{Estación 1 de 3 · Escucha}

\begin{softbox}{milcVerde}{milcGris}{Escena para escuchar}
\textbf{Actividad 1 · IDENTIFICA --- Cuando me acompañaron.} \textbf{Nadie va a leer esta página.} Piensa en \textbf{dos momentos} en que tú la pasaste mal ---una pérdida, una vergüenza, un fracaso, una pelea en la casa--- y responde para cada uno: \textbf{¿qué te dijo la gente?} y \textbf{¿qué te hizo bien y qué te hizo peor?}. Escríbelo con las palabras exactas que recuerdes, incluidas las que te molestaron: «ya se te pasará», «hay gente peor», «no llores», «todo pasa por algo», «yo cuando me pasó\ldots». \textbf{Después, la parte incómoda:} piensa en \textbf{una vez} en que alguien cercano estuvo mal y \textbf{tú no supiste qué hacer}. ¿Qué hiciste? ¿Te acercaste, cambiaste de tema, hiciste un chiste, te alejaste porque no sabías qué decir? \emph{Casi nadie huye por falta de cariño: se huye por no saber qué decir. Y hoy vas a descubrir que casi nunca hay que decir nada.}
\end{softbox}

{\sffamily\bfseries\fontsize{8}{10}\selectfont\textcolor{milcNegro!55}{MARCA CUANDO LO HAYAS HECHO}}
\milccheckrow{Escribiste dos momentos difíciles tuyos con las frases exactas que te dijeron.}
\milccheckrow{Separaste lo que te hizo bien de lo que te hizo peor.}
\milccheckrow{Recordaste una vez en que no supiste acompañar, y qué hiciste en su lugar.}

\begin{infoband}{milcVerde}


\textbf{En tu cuaderno (Actividad 1 · IDENTIFICA).} Título: «Actividad 1. Cuando me acompañaron». \textbf{Formato:} dos situaciones con dos columnas (me hizo bien · me hizo peor) + una situación en que no supiste acompañar. \textbf{Extensión:} media página. \textbf{Sabes que terminaste cuando} tienes al menos \textbf{tres frases textuales} ---dichas por otros--- y no resúmenes. Esta página es tuya: \textbf{nadie más tiene que leerla si tú no quieres}.
\end{infoband}

\puente{Ya los tienes. Ahora la distinción que casi nadie hace: cuatro cosas distintas que en español llamamos igual.}

\milcestacion{milcTurquesa}{Fase 2 · Sistematización: entender la diferencia}

\begin{softbox}{milcTurquesa}{milcGris}{Cuatro claves sobre acompañar sin ponerse por encima}
Cuatro claves para dejar de confundir cuatro cosas que en español se dicen con una sola palabra.
\end{softbox}

\milcpaso{1}{milcTurquesa}{\textbf{Tres cosas distintas, y una trampa.} \textbf{Toma de perspectiva:} entender desde dónde mira el otro ---qué sabe, qué teme, qué está viviendo---; es más de cabeza que de corazón, y se puede entrenar. \textbf{Preocupación empática:} que te importe, que sientas algo por él; es lo que empuja a moverse. \textbf{Malestar personal:} tu propia angustia frente al dolor ajeno (Davis, 1983). \emph{Y esa tercera es la trampa:} si al ver a alguien mal te pones tan mal que terminas siendo tú el que necesita consuelo ---o te vas porque «no lo soportas»---, \textbf{el que está siendo atendido eres tú}. \emph{Reconocerlo no es ser mala persona: es el primer paso para poder estar de verdad.}}
\milcpaso{2}{milcTurquesa}{\textbf{La lástima mira desde arriba; la empatía mira de frente.} La lástima siente pena \emph{por} alguien y, sin querer, lo pone abajo: «pobrecito», «qué tristeza esa familia», «menos mal no me pasó a mí». \textbf{Consuela al que la siente y humilla al que la recibe.} \emph{Se nota en detalles chiquitos:} en el tono, en hablar de la persona en tercera persona estando ella presente, en contar su situación a otros «para que la ayuden». \textbf{La prueba:} si lo que vas a decir o hacer no se lo dirías a alguien que consideras tu igual, es lástima. \emph{Y nadie ---nadie--- quiere ser el pobrecito del salón.}}
\milcpaso{3}{milcTurquesa}{\textbf{Acompañar se hace, no se siente.} Las cantadoras no llegan al velorio a sentir: llegan a \textbf{cantar}, y cantan noches enteras. \emph{Ese es el punto:} acompañar es un conjunto de acciones ---estar, sentarse al lado, llevar algo, quedarse en silencio, volver mañana--- y esas acciones se pueden hacer aunque no sepas qué decir. \textbf{Lo que casi nunca ayuda es lo que solemos hacer:} explicar, dar consejos, buscar el lado bueno, comparar con lo propio, apurar el dolor. \emph{«No sé qué decirte, pero aquí estoy» es una de las mejores frases que existen, y no requiere ninguna sabiduría.}}
\milcpaso{4}{milcTurquesa}{\textbf{Nadie canta solo.} En el alabao hay una voz líder y \textbf{un coro que responde}: nadie sostiene el canto entero por su cuenta. \textbf{Traducción directa:} acompañar a alguien que está muy mal \textbf{no es tarea de una sola persona}, y menos si esa persona tiene trece años. \emph{Si tu amigo está en algo grande ---violencia en la casa, autolesiones, alguien que dice que no quiere vivir---, tu papel es estar y avisarle a un adulto}, no cargarlo solo ni guardar un secreto que te asusta. \textbf{Eso no es traicionar: es hacer coro.}}

\begin{drawbox}{milcTurquesa}{Lo que sí y lo que no, cuando alguien está mal}
\begin{pasoslista}
\item \textbf{Sí:} «no sé qué decirte, pero aquí estoy» · «¿quieres que te acompañe?» · «lo que sientes tiene sentido» · «cuéntame si quieres, y si no, también me quedo».
\item \textbf{No:} «ya se te pasará» · «hay gente peor» · «todo pasa por algo» · «no llores» · «yo cuando me pasó\ldots» (eso convierte su historia en la tuya).
\item \textbf{Acciones que valen más que las frases:} sentarse al lado, guardar silencio, llevarle algo, acompañarlo a un lugar, escribirle al otro día, cumplir lo que prometiste.
\item \textbf{Y siempre:} preguntar antes de contar su situación a alguien más ---salvo si hay peligro, y ahí se avisa a un adulto---.
\end{pasoslista}
\end{drawbox}

\begin{softbox}{milcMostaza}{milcGris}{Errores comunes y su corrección}
\textbf{Convertir su historia en la tuya:} «yo cuando me pasó» corre el foco al que estaba escuchando. \textbf{Buscar el lado bueno:} no lo hay todavía, y exigirlo es una carga más. \textbf{Dar consejos sin que los pidan:} el consejo no pedido suena a que el problema era fácil. \textbf{Hablar de la persona en tercera persona delante de ella:} eso es lástima con público. \textbf{Contarlo a otros «para ayudar»:} sin permiso, eso es exponerlo. \textbf{Cargar solo algo grande:} eso no es lealtad, es quedarse sin coro. \textbf{Desaparecer porque no sabes qué decir:} justamente no hay que decir nada.
\end{softbox}

\begin{infoband}{milcTurquesa}


\textbf{En tu cuaderno (Actividad 2 · EXPLICA).} Título: «Actividad 2. Cuatro cosas distintas». \textbf{Formato:} las cuatro ---toma de perspectiva, preocupación empática, malestar personal, lástima--- con \textbf{un ejemplo tuyo} de cada una, sacado de la Actividad 1. \textbf{Extensión:} media página. \textbf{Sabes que terminaste cuando} identificaste al menos una vez en que lo tuyo fue \textbf{malestar personal} y no acompañamiento ---le pasa a todo el mundo---.
\end{infoband}

\puente{Ya sabes distinguirlas. Es hora de practicar el cambio de silla y de escribir lo que sí y lo que no se dice.}

\milcestacion{milcMagenta}{Fase 3 · Praxis: armar la guía}

\begin{softbox}{milcMagenta}{milcGris}{Cómo se acompaña, en concreto}
Acompañar se aprende haciéndolo. Hoy practicas el cambio de silla, escribes tu guía y ---esto es lo que cuenta--- acompañas a alguien de verdad esta semana.
\end{softbox}

\milcpaso{1}{milcMagenta}{\textbf{Cambiar de silla (12 min, en parejas).} Cada quien toma una situación de conflicto o de dolor en la que estuvo involucrado, y \textbf{la cuenta en voz alta desde el punto de vista del otro}, en primera persona: «yo soy \emph{fulano} y lo que me pasó fue\ldots». \textbf{Regla dura:} no vale terminar explicando por qué el otro estaba equivocado. \emph{Si al final de tu relato el otro sigue quedando mal, no cambiaste de silla: te la llevaste contigo.} Tu pareja escucha y al final dice una sola cosa: si te creyó o no.}
\milcpaso{2}{milcMagenta}{\textbf{Lo que sí y lo que no (8 min).} Copia las dos listas del recuadro y \textbf{agrega a cada una dos frases tuyas}: una que te hayan dicho y te haya servido, y una que te hayan dicho y te haya caído mal ---de tu Actividad 1---. \emph{Tu lista vale más que la del libro, porque son las palabras que se usan de verdad en tu casa y en tu curso.}}
\milcpaso{3}{milcMagenta}{\textbf{El acompañamiento real (esta semana).} Elige a \textbf{una persona} que esté pasando por algo ---no hace falta que sea grave: un duelo, una enfermedad en la casa, una mala racha, alguien que anda solo--- y acompáñala con \textbf{una acción}, no con un discurso: sentarte al lado en el descanso, escribirle, acompañarla a un lugar, llevarle algo, preguntarle al otro día. \textbf{Anótalo después:} qué hiciste, qué dijiste, qué respondió, cómo te sentiste tú. \emph{Y no cuentes lo que te haya contado: eso no se trae a clase.}}
\milcpaso{4}{milcMagenta}{\textbf{Cuidarse uno (5 min).} Escribe dos cosas: \textbf{qué te cuesta a ti} de acompañar ---el silencio, el llanto ajeno, no poder arreglarlo, que se parezca a algo tuyo--- y \textbf{a quién acudirías tú} si acompañar a alguien te deja mal. \emph{Un acompañante sin coro se quema, y ya no le sirve a nadie.}}

\begin{drawbox}{milcMagenta}{Antes de cerrar tu guía}
\begin{checklista}
\item Contaste una situación desde la silla del otro \textbf{sin que quedara mal} al final.
\item Tus dos listas tienen frases propias, no solo las del recuadro.
\item Hiciste ---o tienes fecha para hacer--- un acompañamiento real con una \textbf{acción}, no con un consejo.
\item Escribiste qué te cuesta a ti de acompañar.
\item Sabes a quién acudirías si acompañar a alguien te deja mal.
\item Distingues las situaciones en las que hay que \textbf{avisarle a un adulto} en vez de acompañar solo.
\end{checklista}
\end{drawbox}

\begin{softbox}{milcMagenta}{milcGris}{Plantilla de trabajo}
\textbf{Plantilla de la guía.} \emph{Lo que sí digo:} \_\_\_\_ · \emph{Lo que no digo:} \_\_\_\_ · \emph{Acciones que puedo hacer sin decir nada:} \_\_\_\_ \\[1mm]
\textbf{Mi acompañamiento de la semana.} \emph{A quién:} \_\_\_\_ · \emph{Qué hice:} \_\_\_\_ · \emph{Qué respondió:} \_\_\_\_ · \emph{Cómo me sentí yo:} \_\_\_\_ \\[1mm]
\textbf{Mi cuidado.} \emph{Lo que más me cuesta de acompañar es \_\_\_\_. Si me deja mal, acudo a \_\_\_\_.}
\end{softbox}

\begin{infoband}{milcMagenta}


\textbf{En tu cuaderno (Actividad 3 · APLICA).} Título: «Actividad 3. Mi guía de acompañamiento». \textbf{Formato:} el relato desde la silla del otro + las dos listas con tus frases + el plan del acompañamiento real + tus dos líneas de autocuidado. \textbf{Extensión:} una página. \textbf{Sabes que terminaste cuando} tu guía sirve \textbf{aunque no se te ocurra nada que decir} ---porque está llena de acciones---.
\end{infoband}

\puente{El producto incluye un acompañamiento \textbf{real}. \emph{Esto no se aprende en el cuaderno: se aprende sentándose al lado de alguien.}}

\milcestacion{milcMostaza}{Producto de la sesión}
\textbf{Guía de acompañamiento: cambio de silla, listas propias, un acompañamiento real y el propio cuidado}.
\begin{softbox}{milcMostaza}{milcGris}{Criterios de calidad}
(1) El relato desde el punto de vista del otro es creíble y no termina descalificándolo. (2) Las listas de lo que sí y lo que no incluyen frases propias, tomadas de la experiencia. (3) Hay un acompañamiento real planeado o hecho, basado en una acción concreta. (4) Está escrito qué le cuesta a quien acompaña y a quién acudiría si lo deja mal. (5) Se distinguen las situaciones que exigen avisar a un adulto de las que se acompañan entre pares.
\end{softbox}

\puente{Antes de cerrar, revisa si en tu guía aparece cómo te cuidas tú. Acompañar sin eso dura poco y termina mal.}

\milcestacion{milcVino}{Evaluación liberadora · cinco dimensiones}

\begin{softbox}{milcVino}{milcGris}{Sentido de la evaluación}
Evaluar no es solo revisar una nota. Es reconocer qué aprendí, qué cambió en mí, qué emociones manejé, cómo me relaciono con la ciudad y la comunidad, y qué le dejaré a quien viene detrás.
\end{softbox}

\textbf{Mi reflexión liberadora en cinco dimensiones.} Completa cada frase iniciada:

\small
\begin{tabularx}{\textwidth}{>{\bfseries\RaggedRight\arraybackslash}p{3.4cm}Y>{\RaggedRight\arraybackslash}p{4.6cm}}
\rowcolor{milcVino}\color{white}Dimensión & \color{white}\bfseries Mi respuesta (frame starter) & \color{white}\bfseries Algo que descubrí\\
1. Personal & Lo que más cambió en mí fue \linead{6.5cm} & \linead{4.2cm}\\
2. Emocional & La emoción más fuerte fue \linead{2.5cm} y la regulé \linead{3cm} & \linead{4.2cm}\\
3. Ciudadana & En democracia esto importa porque \linead{6cm} & \linead{4.2cm}\\
4. Local (Cartago) & En mi barrio sería útil para \linead{6.5cm} & \linead{4.2cm}\\
5. Intergeneracional & A mi abuela/abuelo le diría \linead{6.5cm} & \linead{4.2cm}\\
\end{tabularx}
\normalsize

\begin{infoband}{milcVino}
\textbf{Para la próxima sustentación.} Vuelve a esta tabla en seis meses. Verás que las respuestas cambian — no porque lo aprendido cambie, sino porque tú cambias. La evaluación liberadora es una conversación contigo a través del tiempo.
\end{infoband}


\puente{Cerramos con tres voces: escuchar en silencio, lo que puedes hacer cuando no puedes arreglar nada, y qué significa cuidar en un mundo que corre.}

\milcestacion{milcGradeDark}{Triángulo de pensamiento · cierre filosófico}

\begin{softbox}{milcVino}{milcGris}{Dussel · lente del nosotros}
\textit{«Aquel que es capaz de escuchar silenciosamente al Otro como otro es el que puede constituir una comunidad y no una mera sociedad totalizada.»}

{\footnotesize\itshape\color{milcNegro!70}--- Enrique Dussel · Filosofía de la liberación (1977)}

\textbf{«Al Otro como otro»}: no como un problema que resolver, ni como alguien a quien uno mira desde arriba. Esa es exactamente la diferencia entre acompañar y tener lástima. \emph{Y fíjate en la palabra «silenciosamente»}: Dussel no dice escuchar para responder mejor, dice escuchar en silencio ---que es lo que hacen las cantadoras, que no explican nada, solo sostienen el canto---. \textbf{Una comunidad no se hace con gente que da buenos consejos: se hace con gente que se queda.}

\textbf{Cuando alguien me cuenta algo duro, ¿escucho o estoy preparando lo que voy a responder?}
\end{softbox}
\linead{\linewidth}\\[1mm]
\linead{\linewidth}

\begin{softbox}{milcMostaza}{milcGris}{Estoicismo · Marco Aurelio · Meditaciones VII, 47 (c. 175 d.C.) · cuidado interior}
\textit{«Contempla las evoluciones de los astros y piensa al mismo tiempo que evolucionas con ellos; y reflexiona continuamente cómo los elementos cambian unos en otros.»}

{\footnotesize\itshape\color{milcNegro!70}--- Marco Aurelio · Meditaciones VII, 47 (c. 175 d.C.)}

Marco Aurelio pide mirar de lejos \textbf{sin salirse de lo que se mira}: estás adentro, cambias con ello. Eso es justamente la toma de perspectiva bien hecha ---entender desde dónde mira el otro--- \emph{sin} caer en el malestar personal, que es fundirse con su dolor hasta no poder ayudar. \textbf{Ni tan lejos que sea lástima, ni tan cerca que te ahogues.} \emph{Y la segunda parte ---los elementos cambian unos en otros--- sirve de consuelo real: hoy acompañas, y en algún momento del año te va a tocar ser el acompañado.}

\textbf{¿Estoy tan lejos que solo siento lástima, o tan cerca que ya no le sirvo?}
\end{softbox}
\linead{\linewidth}\\[1mm]
\linead{\linewidth}

\begin{softbox}{milcTurquesa}{milcGris}{Floridi · lente de la infoesfera}
\textit{«El florecimiento de las entidades informacionales y de toda la infoesfera debe promoverse preservando, cultivando y enriqueciendo sus propiedades.»}

{\footnotesize\itshape\color{milcNegro!70}--- Luciano Floridi · La ética de la información (2013; trad.)}

Hay una forma muy actual de hacer daño acompañando: \textbf{contar lo que te contaron}. Una historia que alguien te confió en voz baja se puede reenviar a diez personas en un minuto, y ahí deja de pertenecerle a quien la vivió. \emph{Cuidar la infoesfera, en tu curso, es tan concreto como esto:} lo que te dijeron en confianza no se cuenta, no se comenta, no se convierte en tema del grupo. \textbf{La única excepción es el peligro, y ahí no se cuenta: se avisa a un adulto.}

\textbf{¿He contado algo que me confiaron? ¿Qué habría pasado si lo hubieran contado de mí?}
\end{softbox}
\linead{\linewidth}\\[1mm]
\linead{\linewidth}

\begin{infoband}{milcNegro}
\textbf{Nota para el docente.} Las tres son citas textuales y llevan su referencia debajo. Puedes remitir a la obra en clase.
\end{infoband}

\milcestacion{milcVino}{Mi autoevaluación · marca una casilla por fila}

{\small
\setlength{\extrarowheight}{2.2pt}
\begin{tabularx}{\textwidth}{>{\RaggedRight\arraybackslash\bfseries}p{3.0cm}YYY}
\rowcolor{milcVino}\color{white}\bfseries Criterio & \color{white}\bfseries Lo logré & \color{white}\bfseries En proceso & \color{white}\bfseries Necesito apoyo\\[.6mm]
Comprensión del tema & \milcopcion{Explico las ideas con mis palabras.} & \milcopcion{Reconozco las ideas, falta precisión.} & \milcopcion{Necesito repasar lo básico.}\\[.8mm]
\rowcolor{milcGris}Toma de perspectiva & \milcopcion{Puedo contar la situación desde el punto de vista del otro sin descalificarlo.} & \milcopcion{Entiendo al otro pero termino explicando por qué yo tenía razón.} & \milcopcion{Todavía confundo acompañar con dar consejos o con sentir lástima.}\\[.8mm]
Aplicación y producto & \milcopcion{Mi producto cumple los criterios.} & \milcopcion{Está armado, con ajustes pendientes.} & \milcopcion{Aún está en construcción.}\\[.8mm]
\rowcolor{milcGris}Reflexión 5 dimensiones & \milcopcion{Respondí cada una con honestidad.} & \milcopcion{Respondí algunas con detalle.} & \milcopcion{Mis respuestas son muy generales.}\\[.8mm]
Triángulo de pensamiento & \milcopcion{Conecté las 3 voces con mi trabajo.} & \milcopcion{Conecté al menos una.} & \milcopcion{No las relaciono con mi proceso.}\\[.8mm]
\end{tabularx}}

\vspace{3mm}
\textbf{Hoy entendí que:}\\[1mm]
\linead{\linewidth}\\[1mm]
\linead{\linewidth}

\textbf{Mi compromiso para la próxima sesión es:}\\[1mm]
\linead{\linewidth}\\[1mm]
\linead{\linewidth}

\begin{softbox}{milcMostaza}{milcGris}{Cita de despedida · Marco Aurelio}
\textit{``El obstáculo en el camino se vuelve el camino. Lo que estorba al camino se vuelve camino.''}\\[1mm]
Cada error de hoy, bien observado, se convierte en pieza del camino que pisarás mañana.
\end{softbox}

\small
\begin{tabularx}{\textwidth}{>{\bfseries}p{3.6cm}Y}
\rowcolor{milcGradeDark!12}Nombre completo: & \linead{10cm}\\
Fecha: & \linead{4cm} \quad Curso/grupo: \linead{4cm}\\
Mi firma: & \linead{9cm}\\
\end{tabularx}
\normalsize

\begin{infoband}{milcGradeDark}
\textbf{Para la próxima sesión.} Dos cosas. \textbf{Primero, haz el acompañamiento} que planeaste y anota qué hiciste, no qué sentiste ---la lista de acciones es lo que se evalúa---. \emph{No traigas lo que la persona te haya contado: eso se queda con ella.} \textbf{Segundo, busca las cantadoras de tu familia:} pregunta en tu casa \textbf{qué se hace cuando alguien se muere} ---quién llega, quién cocina, quién reza, quién se queda de noche, si se canta, cuántos días dura---. Si tu familia viene del Pacífico, pregunta directamente por los alabaos y los gualíes. Si no, igual va a haber costumbres: la novena, el café, la casa que se llena. \textbf{Trae:} qué hiciste en tu acompañamiento y las costumbres que te contaron. \emph{Vas a descubrir que casi ninguna cultura le pidió nunca a la gente que \emph{sintiera} algo: todas le pidieron que \emph{hiciera} algo, y que lo hiciera acompañada.}
\end{infoband}

\end{document}
